\section{Three split constructions}
\label{sec:splits}

All timing data are now fixed. We use them to define three deterministic
episode updates and their prefix/suffix tests. The main update has the
true monoid action as its linear part. Two larger affine updates encode
the main update and then each other on distinguished basis vectors.
Their graph languages are constructed in reverse order, starting from
an empty skipped class. This section proves semantic correctness of those
tests; Section~\ref{sec:compilation} establishes their expression heights.

There are three different orders to keep separate. Along the word, buffer
$1$ precedes buffer $0$ and the main segment. The total update tables are
defined in the order $g_0,g_1,g_2$. The graph tests are constructed by
stage~1, stage~0, and the main split. Every episode update is fixed on the
entire episode before a split is selected.

Throughout this section, an unavailable inspection window or an unordered
or unavailable product interval makes a test fail.  In particular, every
window belonging to a candidate list is mandatory: a test never removes
an unavailable candidate and then checks uniqueness among those remaining.
The placement in Lemma~\ref{lem:placement} supplies these windows on all
the true uses below, including the windows at false candidate origins.

\subsection{The main update and split}

Fix an episode $e\in E$, with start $s$, origin $t=s+L$, anchor $q$, and
end $b=q+b_{\rm tail}$.  Write
\[
 T_e=T_{s,b},\qquad d_j=T_{B_{j-1},B_j}\quad(1\leq j\leq l),
 \qquad \mathbf d=(d_1,\ldots,d_l).
\]
Recall that, for a residue $\ell\in\mathbb Z/n\mathbb Z$,
\[
 \mathcal E(\ell)
 =\{(\alpha,\lambda):\ell\in I_\alpha+J_\lambda\}
\]
is its clock edge set.  By Lemma~\ref{lem:clock}, this set is either
contained in the diagonal, possibly empty, or consists of one
off-diagonal edge.  With $\ell=q-t\pmod n$, define $g_{0,e}:V\to V$ by
\begin{equation}
\label{eq:main-update}
 g_{0,e}(x)=
 \begin{cases}
 xT_e+\beta(\mathbf d)T_{B_l,b},
   &\mathcal E(\ell)\text{ is contained in the diagonal},\\[2mm]
 xT_e+Y(\lambda)-X(\alpha)T_e,
   &\mathcal E(\ell)=\{(\alpha,\lambda)\},\quad\alpha\ne\lambda.
 \end{cases}
\end{equation}
Here $X(\alpha)$ and $Y(\alpha)$ denote the corresponding fields of a
tag.  In both cases this is an affine map with linear part $x\mapsto xT_e$.
All terms in its definition are fixed by the episode.

We define languages $P_x^{\rm main},Q_y^{\rm main}$ for $x,y\in V$.
In specifying a prefix test, $s$ is the start of its argument; in specifying
a suffix test, $k$ is the start of its argument.  These coordinates are
relative to the argument and do not require an external position origin.

The prefix test $P_x^{\rm main}$ chooses a tag
\[
 \alpha=(j,r,v,d_1,\ldots,d_{j-1},X,Y)
\]
and an integer $0\leq a<n$.  Its argument must end at
\[
 k=t+c_j+\delta_r a,\qquad t=s+L.
\]
It requires
\[
 k-t\pmod n\in I_\alpha,\qquad x=X,\qquad xT_{s,B_0}=v,
\]
and checks the claimed products $d_f=T_{B_{f-1},B_f}$ for $f<j$.
Only boundaries and inspections preceding its cut are used.

The suffix test $Q_y^{\rm main}$ chooses its own tag $\lambda$,
independently of the prefix choice.  First it applies the end guard of
Lemma~\ref{lem:end-guard} to the entire list
\[
 \{k-h:h\in\mathcal H\}.
\]
Thus it requires a common value of $\eta$ at these origins and applies
the exact end rule at a fixed reference origin in the list.  Denote the
resulting anchor by $q$ and the argument end by $b=q+b_{\rm tail}$.
The test requires
\[
 q-k\pmod n\in J_\lambda,\qquad y=Y(\lambda).
\]
Write the selected tag as
$\lambda=(j,r,v,d_1,\ldots,d_{j-1},X,Y)$.
For \emph{every} origin
\[
 \tau=k-c_j-\delta_r a,\qquad 0\leq a<n,
\]
the test computes $B_j(\tau),\ldots,B_l(\tau)$ by the schedule at
$\tau$, and supplies the later products
\[
 d_f^{(\tau)}=T_{B_{f-1}(\tau),B_f(\tau)}\qquad(j<f\leq l).
\]
It checks that the following function of $c\in M$ is constant:
\begin{equation}
\label{eq:main-line-test}
 c\longmapsto
 v\bigl(d_1\cdots d_{j-1}c\,d_{j+1}^{(\tau)}\cdots d_l^{(\tau)}\bigr)
 +\beta\bigl(d_1,\ldots,d_{j-1},c,
                   d_{j+1}^{(\tau)},\ldots,d_l^{(\tau)}\bigr).
\end{equation}
If its constant value is $z_\tau$, it further requires
\[
 z_\tau T_{B_l(\tau),b}=y.
\]
Empty product lists have their usual identity interpretation.  All these
checks are made within the suffix argument.  In particular, the exact end
rule is run only at the reference origin; other origins use its fixed
anchor $q$ and bounded controls.

Two different lists occur in this definition. The end guard uses every
$h\in\mathcal H$, independently of the chosen tag, so it contains the true
origin even when the suffix chooses a wrong role. The line-constancy tests
use all $n$ origins for the suffix's own role $(j,r)$. Both lists are fixed
and retain all their mandatory windows. The first enforces a true end;
the second will use product transport for completeness.

The main split skips episodes whose clock, inner transport, or end
stability may fail. Its residual decoder will use the predicate $A_0$,
but first it must obtain a common anchor from the end guard. For that
decoder we require $A_1$ to be false at every candidate origin, uniformly
over all choices of the cut. These origins are $t+p-p'$ with
$p,p'\in\mathcal P_0$, so their displacements from $t$ form $\Delta_0$.
This requirement determines the second exceptional class. Set
\begin{equation}
\label{eq:residual-domains}
 \begin{split}
 D_0&=\{e\in E:A_0(t,q)\text{ or }A_1(t)\},\\
 D_1&=\{e\in D_0:\text{there exists }\gamma\in\Delta_0
                              \text{ with }A_1(t+\gamma)\},\\
 D_2&=\varnothing.
 \end{split}
\end{equation}
Thus $D_1$ consists of those episodes already in $D_0$ for which the
all-candidate end-stability requirement fails. Since $0\in\Delta_0$,
an episode in $D_0\setminus D_1$ satisfies $A_0(t,q)$ and has $A_1$
false at every $t+\gamma$, $\gamma\in\Delta_0$. Stage~$0$ can therefore
use a common anchor and seek the true origin through $A_0$. Stage~$1$
will instead decode the bounded end-instability predicate in the
definition of $D_1$ before testing that origin's end.

The buffers will implement affine lifts $g_1:U_1\to U_1$ and
$g_2:U_2\to U_2$, defined below. The three planned split applications
are
\begin{center}
\begin{tabular}{@{}llll@{}}
\toprule
Split & Update tested & Episode class & Skipped class \\
\midrule
main & $g_0$ & $E$ & $D_0$ \\
stage $0$ & $g_1$ & $D_0$ & $D_1$ \\
stage $1$ & $g_2$ & $D_1$ & $D_2=\varnothing$ \\
\bottomrule
\end{tabular}
\end{center}
These applications are used in the order stage~1, stage~0, main.

\begin{proposition}[Main split]
\label{prop:main-split}
The languages $P_x^{\rm main},Q_y^{\rm main}$ satisfy the two split
hypotheses of Lemma~\ref{lem:split} for the episode update $g_0$, with
$D=E$ and $D'=D_0$.
\end{proposition}

\begin{proof}
\emph{Soundness for arbitrary independent tags.}
Suppose that a prefix of a word in $E^*$ factors as
$P_x^{\rm main}Q_y^{\rm main}$.  The prefix length is positive and its
cut lies before the earliest possible first episode end, by
Lemma~\ref{lem:placement}.  Hence there is a first episode $e$.
If the prefix chose $k=t+h$ with $h\in\mathcal H$, then the true origin
$t$ belongs to the suffix's end-guard list.  Its diameter is
$\operatorname{diam}\mathcal H<w$.  Lemma~\ref{lem:end-guard} therefore
forces the concatenation to end exactly at the end of $e$, with its true
anchor $q$.  This conclusion applies also to a proposed factor ending
too early or extending into later episodes; it does not presume that the
suffix has already identified the episode correctly.

Let $\alpha$ and $\lambda$ be the independently chosen prefix and suffix
tags.  The two residue tests give
$(\alpha,\lambda)\in\mathcal E(q-t)$.  If this edge set is contained in
the diagonal, then $\alpha=\lambda$.  The true origin is consequently
among the candidates for the selected role $(j,r)$ in the suffix test.
At that candidate, the transmitted vector and preceding products are
correct by the prefix checks, and the later products are correct by the
suffix checks.  Evaluating the constant function
\eqref{eq:main-line-test} at the actual $d_j$ gives the true corrected
vector at $B_l$.  Its transport to $b$ proves $g_{0,e}(x)=y$.
If instead there is one off-diagonal edge, the prefix forces
$x=X(\alpha)$ and the suffix forces $y=Y(\lambda)$.  Formula
\eqref{eq:main-update} again gives $g_{0,e}(x)=y$.  Thus every accepted
split is sound, regardless of which tags were guessed.

\emph{Completeness at one usable scale.}
Let $e\in E\setminus D_0$ and fix $x\in V$.
Then $q-t\pmod n\notin W$, some scale $r$ is usable at $t$, and
$A_1(t)$ fails.  Put $v=xT_{s,B_0}$.  Lemma~\ref{lem:affine-line}
chooses an interval $j$ for which the corrected segment function is
constant on the coordinate line through the true tuple $\mathbf d$.
Choose the tag with this $j,r,v$, the true preceding products, and
\[
 X=x,\qquad Y=g_{0,e}(x).
\]
The full diagonal is available outside $W$.  Since
$\gcd(\delta_r,n)=1$, the cuts $t+c_j+\delta_r a$ for $0\leq a<n$
realize every residue modulo $n$.  One therefore satisfies both clock
tests for this tag.

The end guard succeeds because $A_1(t)$ fails.  Each candidate for the
selected role differs from $t$ by a multiple of $\delta_r$ of magnitude
at most $R_r$.  Lemma~\ref{lem:boundary-transport} leaves all its later
interval products and its $B_l$-to-$b$ product unchanged.  Thus every
candidate in the universal suffix check has the same constant function
and the same final output as the true candidate.  Notice that this uses
only the later products: no claim about the vector at a shifted $B_0$
is needed.  The selected cut gives a split with output $g_{0,e}(x)$.
\end{proof}

\subsection{Total tables and affine lifts}

To handle $D_0$, we encode the entire main update, including its affine
correction, in the linear part of an update on $U_1$. Repeating this
construction on $U_2$ will handle the second exceptional class. The
notation $[v,c]$ always denotes a basis vector of
$U_{i+1}=\F_2^{U_i\times M}$, even when $v=0$; in particular,
$[0,1_M]$ is distinct from the zero vector of $U_{i+1}$.

For each episode, a \emph{total prefix table} at $a_i$ is a family of
maps
\[
 \Phi_i(c):U_i\longrightarrow U_i\qquad(c\in M)
\]
determined by the anchored schedule and the suffix data from $a_i$
onward, such that
$\Phi_i(T_{s,a_i})=g_{i,e}$.
The table must be defined at every $c\in M$, including values that are
not attainable by a prefix of the relevant length.

Here is the initial table explicitly.  Set
\[
 A=T_{a_0,b},\qquad C=T_{B_l,b}.
\]
The main tuple $\mathbf d$ and the clock edge set are determined from
the suffix after $a_0$.  For every $c\in M$ and $v\in U_0=V$, define
\begin{equation}
\label{eq:initial-prefix-table}
 \Phi_0(c)(v)=
 \begin{cases}
 v(cA)+\beta(\mathbf d)C,
       &\mathcal E(q-t)\text{ is contained in the diagonal},\\[2mm]
 v(cA)+Y(\lambda)-X(\alpha)(cA),
       &\mathcal E(q-t)=\{(\alpha,\lambda)\},\quad\alpha\ne\lambda.
 \end{cases}
\end{equation}
All products retain their reading order.  Substitution of
$c=T_{s,a_0}$ gives \eqref{eq:main-update}, while the formula itself
uses no knowledge of that true prefix product.

\begin{lemma}[Total tables and affine lifts]
\label{lem:table-lift}
There are episode updates $g_{1,e}:U_1\to U_1$ and
$g_{2,e}:U_2\to U_2$, affine for every $e\in E$, with the following
properties.  For $i=0,1$, the update $g_{i,e}$ has a total prefix table
$\Phi_i$ at $a_i$.  Define a linear map $H_i:U_{i+1}\to U_{i+1}$ on
basis vectors by
\[
 [v,c]H_i=[\Phi_i(c)(v),1_M].
\]
Writing
$\mathbf d_i=(T_{B_{i,0},B_{i,1}},\ldots,
               T_{B_{i,l_i-1},B_{i,l_i}})$,
the next update is
\begin{equation}
\label{eq:lift-update}
 g_{i+1,e}(z)=
 \left(z\rho_i(T_{s,a_i})
       +\beta_i(\mathbf d_i)\rho_i(T_{B_{i,l_i},a_i})\right)H_i.
\end{equation}
Its linear part $L_{i+1,e}$ satisfies
\begin{equation}
\label{eq:lift-linear-part}
 L_{i+1,e}=\rho_i(T_{s,a_i})H_i,
 \qquad
 [v,1_M]L_{i+1,e}=[g_{i,e}(v),1_M]
 \quad(v\in U_i).
\end{equation}
\end{lemma}

\begin{proof}
The initial total table is \eqref{eq:initial-prefix-table}.
Given any total table $\Phi_i$, prescribing $H_i$ on the basis defines
a unique linear map.  The maps $\Phi_i(c)$ need not be linear as maps
on $U_i$: their values only select the images of basis vectors in the
larger space.  Formula~\eqref{eq:lift-update} is therefore affine, with
the linear part displayed in \eqref{eq:lift-linear-part}.  Moreover,
\[
 [v,1_M]\rho_i(T_{s,a_i})H_i
   =[v,T_{s,a_i}]H_i
   =[\Phi_i(T_{s,a_i})(v),1_M]
   =[g_{i,e}(v),1_M].
\]

The lift has been constructed from a total table. We must still verify
that the next suffix can compute such a table without the earlier prefix.
This is the precise role of the reversed physical buffer order.

It remains to obtain the table $\Phi_1$ needed to construct $g_2$.
The buffer order is $1$ then $0$, so $a_1$ precedes all the segment
data of buffer $0$.  Define, for every $c\in M$ and $z\in U_1$,
\begin{equation}
\label{eq:lift-table}
 \Phi_1(c)(z)=
 \left(z\rho_0(cT_{a_1,a_0})
       +\beta_0(\mathbf d_0)\rho_0(T_{B_{0,l_0},a_0})\right)H_0.
\end{equation}
The actual tuple $\mathbf d_0$, the displayed suffix products, and
the entire map $H_0$ are available from the suffix after $a_1$.
Indeed, $H_0$ is built from the total table $\Phi_0$, whose values
depend only on the later anchored schedule and suffix products; it
does not require the true value of $T_{s,a_0}$. More explicitly, with
$q,t,\mathbf d,A,C$ fixed by those later data, the suffix enumerates
\eqref{eq:initial-prefix-table} for every $c'\in M$ and $v\in U_0$.
Those values give every basis image of $H_0$, including images indexed by
infeasible prefix values. Thus
\eqref{eq:lift-table} is a total table with the required dependence.
Substitution of $c=T_{s,a_1}$ yields \eqref{eq:lift-update} for $g_1$.
The preceding construction then gives $H_1$ and $g_2$.
All these maps are defined from the complete episode independently
of any split or skipped class.
\end{proof}

\begin{lemma}[Recovery on sequences]
\label{lem:sequence-lift}
For any $D\subseteq E$ and $i\in\{0,1\}$, graph tests for $g_{i+1}$
on $D^*$ yield graph tests for $g_i$ on $D^*$ by finite Boolean
operations.  Consequently this recovery does not increase generalized
star height.
\end{lemma}

\begin{proof}
For a sequence $e_1\cdots e_k\in D^*$, the linear part of the affine
composite of $g_{i+1,e_1},\ldots,g_{i+1,e_k}$ is
$L_{i+1,e_1}\cdots L_{i+1,e_k}$, in reading order.
Equation~\eqref{eq:lift-linear-part} shows that the distinguished basis
vectors return to distinguished basis vectors after every episode.
By induction, this product sends $[v,1_M]$ to
\[
 [g_{i,e_k}(\cdots g_{i,e_1}(v)\cdots),1_M].
\]
If $R_{\xi,\zeta}$ are graph tests for the affine composite on
$U_{i+1}$, the graph test for sending $v\in U_i$ to $w\in U_i$ is
therefore
\[
 \bigcup_{f\in U_{i+1}}
 \left(R_{[v,1_M],[w,1_M]+f}\cap R_{0,f}\right),
\]
by Lemma~\ref{lem:affine-recovery}.  The zero in $R_{0,f}$ is the
zero vector of $U_{i+1}$.  This formula applies to the entire sequence,
including the empty sequence, and proves the claim.
\end{proof}

\subsection{Two aligned decoders}

We have obtained affine updates and a recovery identity for entire
sequences. It remains to give the two residual splits. Their algebraic
transmission rule is the same; their origin decoders use different
predicates and a different order of end detection.

For $i=0,1$, we now define $P_x^{(i)},Q_y^{(i)}$ for inputs and outputs
$x,y\in U_{i+1}$.  These tests implement $g_{i+1}$ on $D_i$, with
skipped class $D_{i+1}$.

The prefix test chooses an offset $p\in\mathcal P_i$ and ends at
$k=s+p$.  If the kind assigned to $p$ is
\[
 (j,v,d_1,\ldots,d_{j-1}),
\]
it checks
\[
 v=x\rho_i(T_{s,B_{i,0}}),\qquad
 d_f=T_{B_{i,f-1},B_{i,f}}\quad(1\leq f<j).
\]
The relevant boundaries are the fixed buffer positions relative to
$s$, and all these checks precede the cut.

The suffix lists the candidate origins
\[
 \tau_{p'}=k+L-p'\qquad(p'\in\mathcal P_i).
\]
Its decoding and end checks are as follows.
\begin{enumerate}[label=\textup{(\arabic*)},leftmargin=*]
\item At stage $0$, require $A_1(\tau_{p'})$ to be false for every
candidate.  Apply the end guard to the full candidate list and obtain
its anchor $q$.  Require a unique $p'$ such that
$A_0(\tau_{p'},q)$ holds.  All evaluations of $A_0$ use this same
anchor; they make no additional unbounded end searches.
\item At stage $1$, first require a unique $p'$ such that
\[
 \text{there exists }\gamma\in\Delta_0
                  \text{ with }A_1(\tau_{p'}+\gamma).
\]
This test uses bounded controls and does not involve $q$.
After decoding $p'$, apply the exact end rule at $\tau_{p'}$ to
obtain $q$ and the argument end.
\end{enumerate}
At stage $0$ the candidate list has diameter $h_0<w$, as required
by the end guard.  Stage $1$ does not use that guard and needs no
such bound on $h_1$.

After decoding, set $s'=k-p'$ and use the kind attached to $p'$,
say $(j,v,d_1,\ldots,d_{j-1})$.  All subsequent buffer positions
are scheduled relative to $s'$, and the late controls use
$t'=s'+L$ and the anchor just obtained.  The suffix supplies the
later products $d_f$ for $j<f\leq l_i$ and checks that
\[
 c\longmapsto
 v\rho_i(d_1\cdots d_{j-1}c\,d_{j+1}\cdots d_{l_i})
 +\beta_i(d_1,\ldots,d_{j-1},c,d_{j+1},\ldots,d_{l_i})
 \quad(c\in M)
\]
is constant.  If its value is $z$, the test requires
\[
 \bigl(z\rho_i(T_{B_{i,l_i},a_i})\bigr)H_i=y.
\]
The map $H_i$ is computed from the explicit total tables
\eqref{eq:initial-prefix-table} and \eqref{eq:lift-table}, keeping
the decoded time and anchor data fixed.  These tables call no graph
test on an episode sequence.  The suffix uses only its own data;
if an incorrect hypothesis places required data outside its argument,
the test rejects that argument.

\begin{proposition}[Aligned splits]
\label{prop:aligned-splits}
For $i=0,1$, the languages $P_x^{(i)},Q_y^{(i)}$ satisfy the two
hypotheses of Lemma~\ref{lem:split} for the update $g_{i+1}$, with
$D=D_i$ and $D'=D_{i+1}$.
\end{proposition}

\begin{proof}
\emph{Soundness: decode the true offset and force the true end.}
Consider an accepted concatenation
$P_x^{(i)}Q_y^{(i)}$ at the start of a word in $D_i^*$.
The positive prefix cut lies before every possible first episode
end, so there is a first episode $e\in D_i$.  If the prefix chose
$p$, then its true origin occurs in the suffix list as
\[
 t=s+L=k+L-p=\tau_p.
\]

At stage $0$, the end guard therefore forces exact consumption of
this first episode and supplies its true $q$.  The suffix also
requires $A_1(t)$ to fail.  Since $e\in D_0$, this implies
$A_0(t,q)$.  The true candidate $p$ thus satisfies the decoding
predicate, so uniqueness forces the decoded $p'$ to equal $p$.

At stage $1$, membership of $e$ in $D_1$ directly implies that
$p$ satisfies the decoding predicate.  Every window in that
predicate is mandatory and, on the true episode, lies inside the
episode.  An accepted proposed suffix therefore sees the same
bounded data even if its proposed end was too early or too late.
Uniqueness first forces $p'=p$; the exact end rule at this correctly
decoded origin then forces exact consumption of $e$.

In both stages the prefix and suffix have consequently used the
same kind and the true buffer schedule.  The prefix supplies the
correct vector at $B_{i,0}$ and the preceding interval products;
the suffix supplies the correct later products.  Line constancy
then gives the true corrected vector at $B_{i,l_i}$.  Its subsequent
work action and terminal map $H_i$ give exactly
\eqref{eq:lift-update}, proving $g_{i+1,e}(x)=y$.

\emph{Completeness: count all lost offsets of the needed kind.}
Fix $e\in D_i\setminus D_{i+1}$ and
$x\in U_{i+1}$.  Apply Lemma~\ref{lem:affine-line} to the true
buffer tuple and the vector
\[
 v=x\rho_i(T_{s,B_{i,0}}).
\]
It gives an interval $j$ and thus one required kind
$(j,v,d_1,\ldots,d_{j-1})$.  Every offset of that kind passes the
prefix checks.  There are more than $C_i$ such offsets.

At stage $0$, the condition $e\notin D_1$ makes $A_1$ false at
all $t+\gamma$, $\gamma\in\Delta_0$.  For every cut $p$, all
candidate origins are of this form, so the absence checks and
the stable end guard pass and give the true $q$.  Moreover
$A_0(t,q)$ holds, so the true candidate passes the current
decoding predicate.  If an offset $p$ fails uniqueness, some
$p'\ne p$ satisfies $A_0(t+p-p',q)$.  By
Proposition~\ref{prop:inner-ambiguity}, at most $C_0$ ordered pairs
$(p,p')$ have this property.  In particular at most $C_0$ offsets
$p$ can fail uniqueness.

At stage $1$, true membership in $D_1$ supplies the true candidate
for every cut.  A failure of uniqueness gives an ordered pair
$p\ne p'$ with
\[
 \text{there exists }\gamma\in\Delta_0
             \text{ with }A_1(t+p-p'+\gamma).
\]
Lemma~\ref{lem:end-ambiguity} bounds the number of these pairs,
and hence the number of lost offsets, by $C_1$.

In either stage, a copy of the required kind remains.  At its
correctly decoded origin all later data are available by placement,
and the affine-line check gives the required output
$g_{i+1,e}(x)$.  This proves completeness.
\end{proof}

\subsection{The recovery chain}

The input and output of each split are now established. Figure~\ref{fig:proof-stages}
shows how the graph languages are assembled. A Boolean recovery changes
which update is being tested while leaving the episode domain and the
height bound fixed. A split enlarges that domain and is the only operation
in this chain that introduces an additional star.

\begin{figure}[tbp]
\centering
\begin{tikzpicture}[font=\small,x=1cm,y=1cm,
 box/.style={draw=blue!55!black,fill=blue!5,rounded corners,
 align=center,minimum width=2.6cm,minimum height=0.85cm},
 >=Stealth]
\node[box] (empty) at (0,0) {$g_2$ on $D_2^*$\\height $0$};
\node[box] (g2) at (4,0) {$g_2$ on $D_1^*$\\height $2$};
\node[box] (g1d1) at (4,-1.8) {$g_1$ on $D_1^*$\\height $2$};
\node[box] (g1d0) at (8,-1.8) {$g_1$ on $D_0^*$\\height $3$};
\node[box] (g0d0) at (8,-3.6) {$g_0$ on $D_0^*$\\height $3$};
\node[box] (g0) at (12,-3.6) {$g_0$ on $E^*$\\height $4$};
\draw[->] (empty) -- node[above,align=center,font=\scriptsize]{stage 1\\split} (g2);
\draw[->] (g2) -- node[right,font=\footnotesize]{recover $g_1$} (g1d1);
\draw[->] (g1d1) -- node[above,align=center,font=\scriptsize]{stage 0\\split} (g1d0);
\draw[->] (g1d0) -- node[right,font=\footnotesize]{recover $g_0$} (g0d0);
\draw[->] (g0d0) -- node[above,align=center,font=\scriptsize]{main\\split} (g0);
\end{tikzpicture}

\caption{Construction of the sequence graph languages. Horizontal arrows
apply a split; vertical arrows recover the lower update by finite Boolean
operations. The labels $0,2,3,4$ are the height bounds proved in
Section~\ref{sec:compilation}. Each recovery uses the same word in both
of its affine tests.}
\label{fig:proof-stages}
\end{figure}

\begin{proposition}
\label{prop:recursive-recovery}
Three applications of Lemma~\ref{lem:split}, with intervening finite
Boolean operations, construct tests for the $T$-value on $E^*$ from
the component languages defined above.
\end{proposition}

\begin{proof}
On $D_2^*=\{\eps\}$, the graph tests for $g_2$ are $\{\eps\}$
for equal input and output and empty otherwise.  Apply
Proposition~\ref{prop:aligned-splits} at stage $1$ and
Lemma~\ref{lem:split} to obtain graph tests for $g_2$ on $D_1^*$.
Lemma~\ref{lem:sequence-lift} converts these to graph tests for
$g_1$ on $D_1^*$.

Use those tests as the skipped-class tests in the stage-$0$
application of Proposition~\ref{prop:aligned-splits}.  The split
lemma gives graph tests for $g_1$ on $D_0^*$, and
Lemma~\ref{lem:sequence-lift} recovers graph tests for $g_0$ on
$D_0^*$.  These are the skipped-class tests required by
Proposition~\ref{prop:main-split}, whose application yields graph
tests for $g_0$ on $E^*$.

Finally, Lemma~\ref{lem:affine-recovery} recovers the product of the
linear parts of the maps $g_{0,e}$.  By \eqref{eq:main-update},
these are the true $T_e$ actions on $V$.  Applying their product to
the basis vector indexed by $1_M$ identifies the $T$-value of the
whole episode sequence.  Only the three stated split applications
introduce stars through the split formula; all recovery operations
are finite Boolean operations.  Section~\ref{sec:compilation}
will establish the height bounds for the component languages and
complete the count.
\end{proof}
